On plonge totalement une cellule conductimétrique constituée de deux plaques
parallèles, de surface $S=\qty{1.0}{\square\cm}$, distantes de
$\ell=\qty{1.0}{\cm}$ dans une solution ionique. La tension appliquées entre les
deux électrodes de la cellule est $U=\qty{1.00}{\V}$ et l'intensité électrique
mesurée est $I=\qty{12.0}{\mA}$
\begin{enumerate}
    \item Faire un schéma du montage expérimental.~\textbf{(1pt)}
    \item Déterminer la résistance et la conductance de la portion de solution
          comprise entre les deux électrodes.~\textbf{(1pt)}
    \item Déterminer la conductivité $\sigma$ de la solution.~\textbf{(1pt)}
    \item Quelle serait la valeur de la conductance si on immergeait à moitié
          les électrodes dans la même solution~?~\textbf{(1pt)}
    \item Quelle serait la valeur de la conductance si on divisait par 2 la
          distance séparant les électrodes totalement immergées dans cette même
          solution~?~\textbf{(1pt)}
\end{enumerate}
